\begin{helper}
Let \(\varepsilon_1,\eta\in\mathbb R\) and \(n\in\mathbb N\).  Put
\[
L=\log n,\qquad
K=\noderef{TwoBiteNaturalIndependenceScale}(1+\eta,n),\qquad k=(K:\mathbb R).
\]
If \(0\le k\), \(2\le k\), \(0<\sqrt L\), and
\[
\frac1{\sqrt L}\le\frac{\varepsilon_1}{2},
\]
then
\[
\frac2{\sqrt L}\noderef{RealChooseTwo}(k)
\le \frac{\varepsilon_1}{2}k^2.
\]
\end{helper}
\begin{proof}
By \(\noderef{RealChooseTwoQuadraticBounds}\), the assumptions \(0\le k\)
and \(2\le k\) imply
\[
\noderef{RealChooseTwo}(k)\le \frac{k^2}{2}.
\]
The factor \(2/\sqrt L\) is nonnegative because \(\sqrt L>0\), so multiplying
by this factor preserves the inequality:
\[
\frac2{\sqrt L}\noderef{RealChooseTwo}(k)
\le \frac2{\sqrt L}\cdot\frac{k^2}{2}
=\frac1{\sqrt L}k^2.
\]
Since \(k^2\ge0\), multiplying the threshold
\[
\frac1{\sqrt L}\le\frac{\varepsilon_1}{2}
\]
by \(k^2\) gives
\[
\frac1{\sqrt L}k^2\le\frac{\varepsilon_1}{2}k^2.
\]
Combining the two displayed inequalities proves the claim.
\end{proof}
